\VerdierPhysicalPageAnchor{0112}%
\begingroup
\small
\setlength{\abovedisplayskip}{4pt plus 2pt minus 2pt}%
\setlength{\belowdisplayskip}{4pt plus 2pt minus 2pt}%

\VerdierSemanticAnchor{ch2-1-1-3}
\noindent\textbf{Définition 1.1.3.}\quad
Soient $\mathcal D$ et $\mathcal D'$ deux catégories triangulées. Un
\emph{foncteur exact} de $\mathcal D$ dans $\mathcal D'$ est un
$\mathbb Z$-foncteur (non nécessairement strict) (chap. I, 1.2.4)
$(F,p)\colon\mathcal D\longrightarrow\mathcal D'$ tel que, pour tout
triangle distingué :

\VerdierSemanticAnchor{ch2-1-1-3-source-triangle}
\[
  X\xrightarrow{u}Y\xrightarrow{v}Z\xrightarrow{w}X[1]
\]
de $\mathcal D$, le triangle :

\VerdierSemanticAnchor{ch2-1-1-3-image-triangle}
\[
  FX\xrightarrow{Fu}FY\xrightarrow{Fv}FZ
  \xrightarrow{p(X)Fw}(FX)[1]
\]
soit distingué (1.1.4.2).

\VerdierSemanticAnchor{ch2-1-1-4}
\medskip
\noindent\textbf{Définition 1.1.4.}\quad
Soient $\mathcal D_i$, $1\leq i\leq n$, et $\mathcal D$, $n+1$ catégories
triangulées. Un \emph{foncteur multi-exact} de $\prod_i\mathcal D_i$ dans
$\mathcal D$ est un $\mathbb Z^n$-foncteur tordu (chap. I, 1.4.4)
$(F,p)\colon\prod_i\mathcal D_i\longrightarrow\mathcal D$ tel que, pour
toute famille d'objets $X_i$ appartenant respectivement aux catégories
$\mathcal D_i$, $i\ne j$, le $\mathbb Z$-foncteur :

\VerdierSemanticAnchor{ch2-1-1-4-partial-functor}
\[
  X\longmapsto F(X_1,\ldots,X,\ldots,X_n)\colon
  \mathcal D_j\longrightarrow\mathcal D
\]
soit un foncteur exact.

\VerdierSemanticAnchor{ch2-1-1-4-1}
\medskip
\noindent\textbf{1.1.4.1.}\quad
Dans la définition (1.1.4), la catégorie $\prod_i\mathcal D_i$ est munie
de sa $\mathbb Z^n$-structure canonique et le groupe $\mathbb Z^n$ opère
sur la catégorie $\mathcal D$ par l'intermédiaire de l'homomorphisme

\VerdierSemanticAnchor{ch2-1-1-4-1-group-action}
\[
  \mathbb Z^n\longrightarrow\mathbb Z,
\]
somme des coordonnées (chap. I, 1.6.6).

\VerdierSemanticAnchor{ch2-1-1-4-2}
\medskip
\noindent\textbf{1.1.4.2.}\quad
Conformément à la règle des signes I (chap. I, 1.6.6), les foncteurs
multi-exacts seront, sauf mention expresse du contraire, tordus par le
cocycle de Koszul relatif à la base canonique $e_1,\ldots,e_n$ de
$\mathbb Z^n$ (chap. I, 1.6.4).

\VerdierSemanticAnchor{ch2-1-1-4-2-trivial-cocycle}
Ceci implique en particulier que les foncteurs exacts de $\mathcal D$
dans $\mathcal D'$ sont tordus par le cocycle trivial.

\VerdierSemanticAnchor{ch2-1-1-5}
\medskip
\noindent\textbf{Définition 1.1.5.}\quad
Soient $\mathcal D$ une catégorie triangulée et $\mathcal A$ une catégorie
abélienne. Un foncteur $H\colon\mathcal D\longrightarrow\mathcal A$ est
appelé \emph{foncteur cohomologique} si pour tout triangle distingué :

\VerdierSemanticAnchor{ch2-1-1-5-source-triangle}
\[
  X\xrightarrow{u}Y\xrightarrow{v}Z\xrightarrow{w}X[1]
\]
le diagramme :

\VerdierSemanticAnchor{ch2-1-1-5-image-sequence}
\[
  HX\xrightarrow{Hu}HY\xrightarrow{Hv}HZ\xrightarrow{Hw}H(X[1])
\]
est une suite exacte de $\mathcal A$.
\endgroup
